% Chapter 5: Implementation
\chapter{IMPLEMENTATION}

\section{Development Environment Setup}

\subsection{Required Software}
\begin{itemize}
\item Node.js v18.0.0+
\item npm v9.0.0+
\item Git v2.30.0+
\item VS Code with TypeScript extension
\end{itemize}

\subsection{Project Initialization}
\begin{lstlisting}[language=bash]
# Create Next.js project
npx create-next-app@latest edusphere-ai --typescript

# Install dependencies
npm install @supabase/supabase-js @supabase/ssr
npm install pdf-lib jspdf
npm install -D tailwindcss postcss autoprefixer
\end{lstlisting}

\section{Technology Stack Implementation}

\subsection{Frontend: Next.js 14 + React 18}
Key features used:
\begin{itemize}
\item App Router for file-based routing
\item Server Components for performance
\item Client Components for interactivity
\item API Routes for backend logic
\end{itemize}

\subsection{Backend: Supabase}
\begin{lstlisting}[language=JavaScript]
// lib/supabase.ts
import { createBrowserClient } from '@supabase/ssr'

export function createClient() {
  return createBrowserClient(
    process.env.NEXT_PUBLIC_SUPABASE_URL!,
    process.env.NEXT_PUBLIC_SUPABASE_ANON_KEY!
  )
}
\end{lstlisting}

\section{Key Module Implementation}

\subsection{Authentication Module}

\begin{lstlisting}[language=TypeScript]
// Login functionality
async function handleLogin(email: string, password: string) {
  const supabase = createClient()
  const { data, error } = await supabase.auth.signInWithPassword({
    email,
    password
  })
  if (error) throw error
  return data
}
\end{lstlisting}

\subsection{Workflow Generation Module}

\begin{lstlisting}[language=TypeScript]
function generateWorkflowSteps(context: DocumentContext) {
  const steps = []
  
  if (context.type === 'notice' || context.type === 'circular') {
    if (context.scope === 'department') {
      steps.push({ role: 'hod', description: 'HOD Approval' })
    } else if (context.scope === 'college') {
      steps.push({ role: 'hod', description: 'HOD Approval' })
      steps.push({ role: 'principal', description: 'Principal Approval' })
    }
  }
  
  return steps
}
\end{lstlisting}

\subsection{Digital Signature Module}

\begin{lstlisting}[language=TypeScript]
import { PDFDocument } from 'pdf-lib'

async function embedSignature(
  pdfBytes: Uint8Array,
  signatureImage: Uint8Array,
  x: number, y: number
) {
  const pdfDoc = await PDFDocument.load(pdfBytes)
  const signatureImg = await pdfDoc.embedPng(signatureImage)
  
  const firstPage = pdfDoc.getPage(0)
  firstPage.drawImage(signatureImg, {
    x, y,
    width: 150,
    height: 50
  })
  
  return await pdfDoc.save()
}
\end{lstlisting}

\subsection{Code Execution Module}

\begin{lstlisting}[language=TypeScript]
async function executeCode(code: string, language: string, testCases: any[]) {
  const results = []
  
  for (const testCase of testCases) {
    const response = await fetch('https://judge0-ce.p.rapidapi.com/submissions', {
      method: 'POST',
      headers: {
        'Content-Type': 'application/json',
        'X-RapidAPI-Key': process.env.JUDGE0_API_KEY
      },
      body: JSON.stringify({
        source_code: code,
        language_id: getLanguageId(language),
        stdin: testCase.input
      })
    })
    
    const submission = await response.json()
    results.push(submission)
  }
  
  return results
}
\end{lstlisting}

\section{Database Implementation}

All database queries use Supabase client with type safety:

\begin{lstlisting}[language=TypeScript]
// Fetch documents
const { data: documents } = await supabase
  .from('documents')
  .select('*')
  .eq('creator_id', userId)
  .order('created_at', { ascending: false })

// Insert approval
await supabase.from('approvals').insert({
  workflow_id: workflowId,
  step_order: 1,
  required_role: 'hod',
  status: 'pending'
})
\end{lstlisting}

\section{UI Component Implementation}

\subsection{Reusable Components}

\textbf{Button Component:}
\begin{lstlisting}[language=TypeScript]
export function Button({ 
  children, 
  variant = 'primary',
  onClick 
}: ButtonProps) {
  const baseStyles = 'px-4 py-2 rounded-lg font-semibold'
  const variantStyles = {
    primary: 'bg-blue-600 text-white hover:bg-blue-700',
    secondary: 'bg-gray-200 text-gray-800 hover:bg-gray-300'
  }
  
  return (
    <button 
      className={\\ \\}
      onClick={onClick}
    >
      {children}
    </button>
  )
}
\end{lstlisting}

\section{API Routes Implementation}

\begin{lstlisting}[language=TypeScript]
// app/api/documents/upload/route.ts
export async function POST(request: Request) {
  const supabase = createServerClient()
  
  // Verify authentication
  const { data: { session } } = await supabase.auth.getSession()
  if (!session) {
    return new Response('Unauthorized', { status: 401 })
  }
  
  // Process upload
  const formData = await request.formData()
  const file = formData.get('file') as File
  
  // Upload to storage
  const { data, error } = await supabase.storage
    .from('documents')
    .upload(\\/\\, file)
    
  if (error) throw error
  
  // Create document record
  await supabase.from('documents').insert({
    title: formData.get('title'),
    file_url: data.path,
    creator_id: userId
  })
  
  return Response.json({ success: true })
}
\end{lstlisting}

\section{Challenges and Solutions}

\subsection{Challenge 1: PDF Signature Placement}
\textbf{Problem:} Dynamically placing signatures at user-specified coordinates.\\
\textbf{Solution:} Implemented PDFSignaturePlacer component with click-to-place functionality.

\subsection{Challenge 2: Code Execution Security}
\textbf{Problem:} Ensuring safe execution of user-submitted code.\\
\textbf{Solution:} Used Judge0 sandboxed execution environment with time and memory limits.

\subsection{Challenge 3: Role-Based Access Control}
\textbf{Problem:} Ensuring users only see relevant data.\\
\textbf{Solution:} Implemented row-level security in Supabase and query filtering in API routes.

\section{Development Best Practices}

\begin{itemize}
\item TypeScript strict mode for type safety
\item ESLint for code quality
\item Modular component architecture
\item Consistent naming conventions
\item Comprehensive error handling
\item Git commits with descriptive messages
\item Code reviews before merging
\end{itemize}

\section{Conclusion}

Implementation chapter demonstrates practical application of design decisions. Modern technology stack enables rapid development with robust features. Modular architecture ensures maintainability and scalability. Best practices followed throughout development ensure code quality and reliability.

\newpage
